\documentclass[11pt,a4paper]{article}

\usepackage[utf8]{inputenc}
\usepackage{amsmath, amssymb}
\usepackage{geometry}
\usepackage{hyperref}
\usepackage{listings}
\usepackage{xcolor}
\usepackage{booktabs}

\geometry{margin=2.5cm}

\hypersetup{
    colorlinks=true,
    linkcolor=blue,
    urlcolor=blue,
}

\lstset{
    basicstyle=\ttfamily\small,
    backgroundcolor=\color{gray!10},
    frame=single,
    breaklines=true,
    columns=fullflexible,
    xleftmargin=0.5cm,
    xrightmargin=0.5cm,
}

\title{Bayesian Comparison of TRANSP Runs against \\ Neutron Rate and TH/BT Measurements}
\author{}
\date{}

\begin{document}
\maketitle

\section{Overview}

The script \texttt{bayes\_compare.py} ranks a set of TRANSP runs (different
tritium profiles, same pulse) by how well their predictions match two
experimental observables:

\begin{itemize}
    \item the total neutron rate measured by the fission chamber
          (\textsc{TIN/RNT});
    \item the TH/BT ratio (thermal over beam-target neutron rate) inferred
          from neutron spectroscopy.
\end{itemize}

Each run is treated as a discrete hypothesis $M_i$ and assigned a posterior
probability under a uniform prior. The output is a posterior distribution
over the supplied runs, plus the most likely one.

\section{Bayesian Framework}

\subsection{Posterior, prior, likelihood}

Bayes' theorem gives the posterior probability of run $M_i$ given the data
$\mathbf{d}$:
\begin{equation}
    P(M_i \mid \mathbf{d}) \;=\;
    \frac{P(\mathbf{d} \mid M_i)\, P(M_i)}{P(\mathbf{d})}.
\end{equation}
With a flat prior over the $N$ supplied runs,
$P(M_i) = 1/N$, and the marginal likelihood $P(\mathbf{d})$ is just a
normalisation constant. The posterior is therefore proportional to the
likelihood:
\begin{equation}
    P(M_i \mid \mathbf{d}) \;\propto\; \mathcal{L}(\mathbf{d} \mid M_i).
\end{equation}

\subsection{Gaussian likelihood}

The two measurements are assumed independent and Gaussian-distributed
around the model predictions:
\begin{equation}
    \mathcal{L}(\mathbf{d} \mid M_i)
      \;=\;
      \prod_{j \in \{\text{RNT},\,\text{TH/BT}\}}
      \frac{1}{\sqrt{2\pi \sigma_j^2}}
      \exp\!\left[
          -\frac{1}{2} \left( \frac{d_j - m_{i,j}}{\sigma_j} \right)^{\!2}
      \right].
\end{equation}
Taking the logarithm and dropping the additive constant that is common to
all $M_i$:
\begin{equation}
    \log \mathcal{L}_i \;=\; -\tfrac{1}{2}\,\chi_i^2,
    \qquad
    \chi_i^2 \;=\;
        \left( \frac{\Delta\text{NEUTT}_i}{\sigma_{\text{RNT}}} \right)^{\!2}
        +
        \left( \frac{\Delta(\text{TH/BT})_i}{\sigma_{\text{TH/BT}}} \right)^{\!2},
    \label{eq:chi2}
\end{equation}
where $\Delta X_i = \langle X \rangle_i^{\text{TRANSP}} - X^{\text{meas}}$.

\subsection{Posterior probabilities}

Combining the equal-prior assumption with~\eqref{eq:chi2}:
\begin{equation}
    P(M_i \mid \mathbf{d})
      \;=\;
      \frac{\exp(\log \mathcal{L}_i)}
           {\sum_{k=1}^{N} \exp(\log \mathcal{L}_k)}.
    \label{eq:posterior}
\end{equation}
This is a softmax over the log-likelihoods. For numerical stability the
script subtracts $\max_k \log \mathcal{L}_k$ from every $\log \mathcal{L}_i$
before exponentiating; this leaves the ratio unchanged but keeps the
exponents bounded.

\section{Implementation Details}

\subsection{Time averaging}

Predictions and measurement are compared as time-averages over a
user-supplied window $[t_1, t_2]$ that should fall on a quasi-stationary
phase of the discharge. For a TRANSP signal $X(t)$ on a non-uniform
time grid the average is
\begin{equation}
    \langle X \rangle
      \;=\;
      \frac{\sum_{t_n \in [t_1, t_2]} X(t_n)\, \Delta t_n}
           {\sum_{t_n \in [t_1, t_2]} \Delta t_n},
\end{equation}
using the time weights $\Delta t_n$ already provided by the
\texttt{Transp} class (a trapezoidal-style quadrature). For the RNT signal
the time grid is regular, so a plain mean is used.

\subsection{TH/BT ratio}

For each run the predicted TH/BT is computed as the ratio of
time-averages,
\begin{equation}
    (\text{TH/BT})_i \;=\;
        \frac{\langle \text{NEUTX} \rangle_i}
             {\langle \text{BTNTS} \rangle_i},
\end{equation}
rather than the time-average of the instantaneous ratio. The two are
identical when both signals are constant; the ratio of averages is more
robust when either denominator approaches zero at the boundaries of the
window, and it more closely matches what an integrating
spectroscopy diagnostic measures.

\subsection{Uncertainties}

\begin{itemize}
    \item \textbf{RNT}: relative uncertainty,
          $\sigma_{\text{RNT}} = \alpha_{\text{RNT}} \cdot \langle \text{RNT} \rangle$,
          with default $\alpha_{\text{RNT}} = 0.10$.
    \item \textbf{TH/BT}: absolute uncertainty $\sigma_{\text{TH/BT}}$,
          default $0.06$ (consistent with a $\pm 10\%$ spectroscopy
          measurement around $\text{TH/BT} \simeq 0.6$).
\end{itemize}

\section{Usage}

\subsection{Command line}

\begin{lstlisting}[language=bash]
./bayes_compare.py 104614 M26 M27 M28 \
    -tw 51.5 52.5 \
    --thbt 0.60 \
    --thbt_sigma 0.06 \
    --rnt_unc 0.10
\end{lstlisting}

\subsection{Arguments}

\begin{tabular}{@{}lll@{}}
\toprule
Argument                & Type     & Description \\
\midrule
\texttt{pulse}              & int      & JET pulse number \\
\texttt{runs}               & list     & TRANSP run IDs to compare \\
\texttt{-tw, --time\_window} & 2 floats & Wall-time window $[t_1, t_2]$ in s \\
\texttt{--thbt}              & float    & Measured TH/BT ratio \\
\texttt{--thbt\_sigma}       & float    & Absolute $1\sigma$ on TH/BT (default 0.06) \\
\texttt{--rnt\_unc}          & float    & Relative $1\sigma$ on RNT (default 0.10) \\
\bottomrule
\end{tabular}

\subsection{Sample output}

\begin{lstlisting}
Pulse 104614, time window [51.500, 52.500] s

Measurement:
  <RNT>  = 1.234e+18 n/s   (sigma = 1.234e+17, 10.0%)
  TH/BT  = 0.600           (sigma = 0.060)

Run     NEUTT [n/s]     TH/BT     chi2(RNT)   chi2(TH/BT)   log L         P(M|d)
--------------------------------------------------------------------------------
M26     1.20e+18        0.55      0.072       0.694         -0.383        29.84%
M27     1.30e+18        0.62      0.291       0.111         -0.201        35.78%
M28     1.10e+18        0.58      0.470       0.111         -0.290        34.38%

Most likely: M27  (P = 35.8%)
\end{lstlisting}

The two $\chi^2$ contributions are printed separately so that the user can
see which measurement is driving the discrimination.

\section{Caveats}

\subsection{Independence assumption}

The likelihood is factorised under the assumption that the RNT and TH/BT
measurements are statistically independent. This is reasonable because
they come from physically distinct diagnostics (fission chamber vs.\
neutron spectrometer), but if the two diagnostics share a common
calibration source any correlated systematic would not be captured by the
diagonal covariance used here. A correlated systematic could be added by
replacing the sum in~\eqref{eq:chi2} with $\boldsymbol{\Delta}^\top
\Sigma^{-1} \boldsymbol{\Delta}$, where $\Sigma$ is the full
$2 \times 2$ covariance matrix.

\subsection{Choice of time window}

The result depends critically on the window $[t_1, t_2]$. The window
should be chosen to lie entirely within a stationary phase of the
discharge (typically the flat-top), where the model output and the
measurement can be averaged meaningfully. Including transient phases
inflates the apparent disagreement between the runs and the data and
biases the posterior. If different runs disagree most strongly during a
specific sub-interval, narrowing the window to that interval increases
the discriminating power of the analysis.

\subsection{Discrete prior}

The posterior is defined only over the discrete set of runs that the
user supplies. Two runs receiving similar posterior probabilities does
\emph{not} mean they are equally good descriptions of reality --- it
may instead mean that the truth lies somewhere \emph{between} their
respective tritium profiles, in a region of parameter space not sampled
by any of the runs. The Bayesian probabilities are conditional on the
hypothesis that the true model is one of the runs provided.

\subsection{Statistical vs.\ systematic uncertainty}

The default uncertainties are intended to capture the dominant
systematic component of each diagnostic (calibration, geometric
factors). Statistical uncertainty in the time-averaged measurement is
typically smaller and is not added in quadrature here. If the
diagnostic time series is noisy and the averaging window short, the
statistical contribution should be evaluated separately and folded in.

\subsection{Sensitivity to single time slices}

Although the script averages over a window, the underlying model
predictions are still drawn from a single TRANSP simulation per run.
Differences between runs that are smaller than typical TRANSP-internal
modelling uncertainties (e.g.\ transport coefficients, recycling, fast
ion physics) cannot be reliably resolved by this analysis, even if the
posterior assigns one run a notably higher probability than another.

\end{document}
